\hnsection{THE THEORY OF LIE ALGEBRAS}

Having once acquired the correspondence between transformation groups and Lie algebras, the theory took a considerably more algebraic turn and became centred on a deep study of Lie algebras.†

A first, short period, from 1888 to 1894, marked by the works of Engel and his pupil Umlauf and especially Killing and E. Cartan, achieved a series of spectacular results on complex Lie algebras. We have seen earlier that the notion of solvable Lie algebra was due to Lie himself, who had shown (in the complex case) the theorem on the reduction of solvable linear Lie algebras to triangular form ({[}4{]}, vol.~1, p.~270). \(^{\ddagger}\) Killing observed {[}11{]} that there exists in a Lie algebra a largest solvable ideal (which we now call the radical) and that the quotient of the Lie algebra by its radical has zero radical; he called Lie algebras with zero radical semi-simple and proved that they are products of simple algebras (the latter notion having already been introduced by Lie, who had proved the simplicity of the ``classical'' algebras ({[}4{]}, vol.~3, p.~682)).

On the other hand, Killing introduced in a Lie algebra the characteristic equation \(\det(\mathrm{ad}(x)-\omega.1)=0\) , already encountered by Lie when studying 2-dimensional Lie subalgebras containing a given element of a Lie algebra. We shall return in other Historical Notes to this Book to the analysis of the methods by which Killing, whilst engaged on a penetrating study of the properties of the roots of the ``generic'' characteristic equation for a semi-simple algebra, achieved the most remarkable of his results, the complete determination of the (complex) simple Lie algebras.§

Killing proved that the derived algebra of a solvable algebra is ``of rank 0'' (which means that ad x is nilpotent for every element x of the algebra). Almost immediately, Engel showed that algebras ``of rank 0'' are solvable (this statement is essentially what we called Engel's Theorem in Chapter I, § 4, no. 2). In his thesis, E. Cartan introduced, on the other hand, what we now call the ``Killing form'' and established the two fundamental criteria which characterize by means of this form solvable Lie algebras and semi simple Lie algebras.

Killing had affirmed ({[}11{]}, IV) that the derived algebra of a Lie algebra is the sum of a semi-simple algebra and its radical, which is nilpotent, but his proof was incomplete. A little later, E. Cartan announced without proof ({[}12{]}, vol.~I₁, p.~104) that more generally every Lie algebra is the sum of its radical and a semi-simple subalgebra; the only result in this direction established indisputably at this period was a theorem of Engel affirming the existence, in every non-solvable Lie algebra, of a simple Lie subalgebra of dimension 3. The first published proof (for complex Lie algebras) of Cartan's statement is due to E. E. Levi {[}18{]}; another proof (equally valid in the real case) was given by J. H. C. Whitehead in 1936 {[}26 a{]}. In 1942, A. Malcev completed this result by the uniqueness theorem for ``Levi sections'' up to conjugation.

From his very first works, Lie was concerned with the problem of the isomorphism between any Lie algebra and a linear Lie algebra. He had believed that he had solved it affirmatively by considering the adjoint representation (and deduced from it a proof of his ``third theorem'') \([3\ d]\) ; he realized that his proof was correct only for Lie algebras with zero centre, gave another erroneous proof thereof ( \([3\ f]\) , § 3, no. 9) and then recognized ( \([3\ g]\) , p.~231) that the question remained open. In fact, it remained so for a long time and was only finally resolved affirmatively in 1935 by Ado {[}27{]}. On the other hand, Lie had essentially posed the problem of determining linear representations of simple Lie algebras of minimal dimension and had solved it for the classical groups; in his Thesis, Cartan also solved this problem for the exceptional simple algebras \(^{\dagger}\) ; the methods he used for this were to be generalized by him twenty years later to obtain all the irreducible representations of real or complex simple Lie algebras.

The property of complete reducibility of a linear representation seems to have been encountered for the first time (in a geometric form) by Study. In an unpublished manuscript, but quoted in {[}4{]}, vol.~3, pp.~785--788, he proved this property for linear representations of the Lie algebra \(\mathbf{SL}(2,\mathbf{C})\) and obtained partial results for \(\mathbf{SL}(3,\mathbf{C})\) and \(\mathbf{SL}(4,\mathbf{C})\) . Lie and Engel conjectured on this occasion that the complete reducibility theorem was true for \(\mathbf{SL}(n,\mathbf{C})\) for all \(n\) . The complete reducibility of linear representations of semi-simple Lie algebras was established by H. Weyl in 1925† by a global type of argument (see later). The first algebraic proof was obtained in 1935 by Casimir and van der Waerden {[}32{]}; other algebraic proofs have since been given by R. Brauer {[}31{]} (this is the one we have reproduced) and J. H. C. Whitehead {[}26 b{]}.

Finally, in the course of his research on the exponential mapping (cf.~infra), H. Poincaré ({[}14{]}, vol.~3) considered the associative algebra of differential operators of all orders, generated by the operators of a Lie algebra; he showed essentially that, if \((\mathbf{X}_i)_{1\leqslant i\leqslant n}\) is a basis of the Lie algebra, the associative algebra generated by the \(\mathbf{X}_i\) has as basis certain symmetric functions of the \(\mathbf{X}_i\) (sum of the non-commutative ``monomials'' derived from a given monomial by all permutations of factors). The essence of his proof was algebraic in nature and allowed him to obtain the enveloping algebra structure which we define abstractly in Chapter I. Analogous proofs were given in 1937 by G. Birkhoff {[}29 b{]} and E. Witt {[}30{]}.

Most of the results cited above are limited to real or complex Lie algebras which alone correspond to Lie groups in the usual sense. The study of Lie algebras over a field other than \(\mathbf{R}\) or \(\mathbf{C}\) was embarked upon by Jacobson {[}28 a{]} who showed that the great majority of the classical results (i.e.~those of Chapter I) remain true for any field of characteristic zero.

